\documentclass[12pt]{article}
\usepackage{../../includes/assignment}
\usepackage{../../includes/math}

\begin{document}
\begin{center}
  {\Huge\bf Midterm - Spring 2026}

  \smallskip
  {\large\it ECON 5753 — University of Arkansas}
\end{center}

\medskip
\begin{enumerate}
  \item Consider a researcher drawing a random sample of people from the US and calculating the sample mean of a variable.

    \begin{enumerate}
      \item Describe what the `sampling distribution' of the sample mean is.

      \item Suppose the sampling distribution of $\bar{X}_n$ is approximately normally distributed, i.e. $\bar{X}_n \approx N(\mu, \sigma^2/n)$.
        Explain why this result is useful for quantifying uncertainty about your sample mean.

      \item Given that $\bar{X}_n \approx N(\mu, \sigma^2/n)$, can I say that the sample mean is (i) an unbiased estimator of the population mean and (ii) consistent estimator of the population mean?
        Explain

      \item Say the varaible we are recording is a person's income and we observe 25 individuals in our sample.
        Should I be concerned about using the normal approximation for inference of $\bar{X}_n$?
        Why or why not?
    \end{enumerate}

    \bigskip
  \item For each of the following research questions, classify whether the goal is \textbf{forecasting} or \textbf{inference}.

    \begin{enumerate}
      \item A real estate company wants to predict the sale price of a house based on its square footage, number of bedrooms, and location.

      \item A city planner wants to know how much additional traffic congestion is caused by building a new shopping center on a major road.

      \item A sports team wants to predict which college basketball players will have successful NBA careers based on their college statistics.

      \item A public health researcher wants to understand how much of the variation in life expectancy across counties can be explained by differences in median income.
    \end{enumerate}

    \newpage
  \item You are analyzing data on NBA players (from the video game NBA 2K25). You regress points scored per game on height (in inches).
    Results with indicators for the player's position and without are shown below.
    The omitted category is center (C).

    \emph{Context:} In the NBA, taller positions (`C' center and `PF' power forward) tend to play near the basket and score fewer points than shorter positions (`PG' point guard, `SG' shooting guard, and `SF' small forward) who handle the ball more and shoot from distance.

    \begin{enumerate}
      \item Looking at the results with position controls, interpret the marginal predictive effect of the coefficient on height.

      \item Why do you think the coefficient on height got larger after controlling for position played?

      \item What are two ways you could make the regression model more flexible?

      \item Why might we want to cluster our standard errors at the team level?
    \end{enumerate}

    \begin{codeblock}[{}]
Dependent Var.:         points           points

Constant        -1.716 (9.238)   -9.296 (10.10)
height_in        1.929 (1.526)   2.993. (1.629)
position_1 = PF                  0.9084 (1.053)
position_1 = PG                  2.346* (1.160)
position_1 = SF                  0.6490 (1.016)
position_1 = SG                  1.733. (1.020)
_______________ ______________   ______________
S.E. type       Heterosk.-rob.   Heterosk.-rob.
Observations               432              432
R2                     0.00434          0.01686
---
Signif. codes: 0 '***' 0.001 '**' 0.01 '*' 0.05 '.' 0.1 ' ' 1
    \end{codeblock}

    \newpage
  \item Consider the following scenario: You are interested in studying the relationship
    between college attendance and wages. Let $w_i$ denote hourly wages and $D_i$ be an
    indicator for college attendance ($D_i=1$ if attended, $D_i=0$ otherwise).

    \begin{enumerate}
      \item Suppose you have data on college attendance ($D_i$), gender ($F_i=1$ indicates female), and wages ($Y_i$). 
      Explain how you would estimate the conditional expectation function using the subgroup averages approach.
      Write out all the subgroup averages you would need to compute.

      \item Suppose you also observe each person's race/ethnicity, which has five
        distinct categories (e.g., White, Black, Hispanic, Asian, Other). How many
        subgroup averages would you now need to estimate?
        Explain why this approach becomes impractical as you start adding more variables?
    \end{enumerate}

    \bigskip
  \item Below we present results from a regression using data from the National Longitudinal Survey of Young Men.
    We run a regression of the \emph{logged wages} earned by the worker on the number of years of education they completed.
    In a second regression, we control for the number of years of education that the mother and father completed.
    \begin{enumerate}
      \item Interpret the coefficient on \texttt{educ} (measured in years of schooling) in the second column.
        Discuss the statistical significance of the estimate.
        Note I intentionally removed significance stars, so do not rely on them to answer.

      \item Why do you think the coefficient changed the way it did after adding the two control variables
    \end{enumerate}

    \begin{codeblock}[{}]
Dependent Var.:            ln_wage              ln_wage
Constant           0.9651 (0.0384)      0.9966 (0.0493)
educ               0.0522 (0.0029)      0.0393 (0.0039)
father_educ                             0.0036 (0.0034)
mother_educ                             0.0107 (0.0040)
_______________ __________________   __________________
S.E. type       Heteroskedas.-rob.   Heteroskedas.-rob.
Observations                 3,078                2,271
    \end{codeblock}

\end{enumerate}
\end{document}
